\section*{Chapter \ref{ch:g12:climate-change} --- Climate Change: Observations, Causes and Projections}

\begin{solution}{exo:g12:climate-change:1}
The change it causes in the net flow of energy entering the climate system
at the top of the atmosphere, before the surface temperature responds;
measured in \unit{W/m^2}. Positive: the planet gains energy and warms;
negative: it loses energy and cools.
\end{solution}

\begin{solution}{exo:g12:climate-change:2}
Highest in May, lowest in September. In spring and summer the vegetation
of the Northern Hemisphere, where most land plants grow, takes up more
carbon dioxide by photosynthesis than plants and soils release; in autumn
and winter decay and respiration dominate and the concentration rises again.
\end{solution}

\begin{solution}{exo:g12:climate-change:3}
\emph{Very likely}: a probability of 90 to 100\,\%; \emph{likely}: 66 to
100\,\%. $[0.95 \text{ to } 1.20]$ is a very likely range: a 90\,\% chance that
the true value lies between 0.95 and 1.20.
\end{solution}

\begin{solution}{exo:g12:climate-change:4}
91\,\%. Water expands when it warms: a warmer ocean occupies more volume, so
its surface rises even if no water is added (melting land ice adds water on
top of this).
\end{solution}

\begin{solution}{exo:g12:climate-change:5}
Water vapour: warmer air holds more water vapour, a greenhouse gas, which adds
to the warming. Snow and ice: melting uncovers darker surfaces, which absorb
more sunlight. The negative response: a warmer planet radiates more infrared
to space, which brings it back to balance.
\end{solution}

\begin{solution}{exo:g12:climate-change:6}
$427.35 - 315.98 = \qty{111.37}{ppm}$; $111.37/315.98 \approx 0.352$, an
increase of about 35\,\%. Over 66 years: $111.37/66 \approx
\qty{1.7}{ppm}$ per year on average.
\end{solution}

\begin{solution}{exo:g12:climate-change:7}
$F = 3.93 \times \ln(427.35/278)/\ln 2 = 3.93 \times 0.4300/0.6931 \approx
\qty{2.44}{W/m^2}$.
\end{solution}

\begin{solution}{exo:g12:climate-change:8}
$3.7 \times 30 = \qty{111}{mm}$, about \qty{11}{cm}. $3.7/1.3 \approx 2.8$:
almost three times faster.
\end{solution}

\begin{solution}{exo:g12:climate-change:9}
$\Delta T = 2.72/1.16 \approx \qty{2.3}{\degreeCelsius}$. The climate is not
at equilibrium: the ocean, which takes up about 90\,\% of the extra energy,
warms slowly, so the surface has not yet caught up with the \omterm{def:g12:climate-change:forcing}{forcing} (the
system still gains about \qty{0.8}{W/m^2}).
\end{solution}

\begin{solution}{exo:g12:climate-change:10}
The Sun and volcanoes explain at most \qty{0.1}{\degreeCelsius} of a
\qty{1.06}{\degreeCelsius} warming: they cannot be its cause. The other
natural cause is the internal variability of the climate (its natural
fluctuations, such as exchanges of heat between ocean and atmosphere); the
assessment puts its contribution at $-0.2$ to $+\qty{0.2}{\degreeCelsius}$,
also far too small.
\end{solution}

\begin{solution}{exo:g12:climate-change:11}
Very likely below \qty{2}{\degreeCelsius}: only SSP1-1.9 (very likely range
1.0 to 1.8). Very likely above: SSP2-4.5 (2.1 to 3.5), SSP3-7.0 (2.8 to 4.6)
and SSP5-8.5 (3.3 to 5.7). For SSP1-2.6 (1.3 to 2.4) neither is very likely.
\end{solution}

\begin{solution}{exo:g12:climate-change:12}
The aerosol \omterm{def:g12:climate-change:forcing}{forcing} is $-0.22 - 0.84 = \qty{-1.06}{W/m^2}$; without it the
total becomes $2.72 + 1.06 = \qty{3.78}{W/m^2}$. The equilibrium warming
increases by $1.06/1.16 \approx \qty{0.9}{\degreeCelsius}$.
\end{solution}

\begin{solution}{exo:g12:climate-change:13}
Since $\ln(2C) - \ln C = \ln 2$ whatever $C$,
\[
  F(556) - F(278) = \frac{3.93\,\ln 2}{\ln 2} = \qty{3.93}{W/m^2}, \qquad
  F(1112) - F(556) = \frac{3.93\,\ln 2}{\ln 2} = \qty{3.93}{W/m^2}:
\]
the same \omterm{def:g12:climate-change:forcing}{forcing} for each doubling. Total at \qty{1112}{ppm}: $2 \times 3.93 =
\qty{7.86}{W/m^2}$.
\end{solution}

\begin{solution}{exo:g12:climate-change:14}
$\dfrac{\dd F}{\dd C} = \dfrac{F_{2\times}}{C \ln 2}$. At \qty{280}{ppm}:
$3.93/(280 \times 0.6931) \approx \qty{0.0202}{W/m^2}$ per ppm; at
\qty{420}{ppm}: $\approx \qty{0.0135}{W/m^2}$ per ppm. Each extra ppm adds
less \omterm{def:g12:climate-change:forcing}{forcing} as the concentration grows (1.5 times less at 420 than at 280),
but the concentration keeps rising, so the \omterm{def:g12:climate-change:forcing}{forcing} still grows.
\end{solution}

\begin{solution}{exo:g12:climate-change:15}
$3.93/3.22 \approx \qty{1.2}{\degreeCelsius}$ without \omterm{def:g12:climate-change:feedback}{feedbacks};
$3.93/1.16 \approx \qty{3.4}{\degreeCelsius}$ with them. Amplification
factor: $3.22/1.16 \approx 2.8$.
\end{solution}

\begin{solution}{pb:g12:climate-change:1}
\textbf{1.} $3.37 + 3.51 = \qty{6.88}{ppm}$, about \qty{7}{ppm}.

\textbf{2.} From spring to late summer, the plants of the Northern
Hemisphere draw carbon dioxide out of the air by photosynthesis faster than
respiration and decay return it: the concentration falls until September.
In autumn and winter the balance reverses, and the concentration climbs
back to its May maximum.

\textbf{3.} The Southern Hemisphere has much less land, and so much less
vegetation, than the Northern; most of its surface is ocean.

\textbf{4.} High up, it lies above the layer of polluted air near the
ground; far from continents and cities, the air it samples is well mixed and
represents the whole atmosphere, not a local source.

\textbf{5.} $(325.68 - 316.91)/10 \approx \qty{0.88}{ppm}$ per year.

\textbf{6.} $(427.35 - 401.01)/10 \approx \qty{2.63}{ppm}$ per year.

\textbf{7.} $2.634/0.877 \approx 3.0$: the growth rate has tripled.

\textbf{8.} $(427.35 - 278)/278 \approx 0.537$: about 54\,\%.

\textbf{9.} $427.35 - 300 \approx \qty{127}{ppm}$.

\textbf{10.} $F = 3.93 \times \ln(427.35/278)/\ln 2 \approx
\qty{2.44}{W/m^2}$.

\textbf{11.} $F(316.91) = 3.93 \times \ln(316.91/278)/\ln 2 \approx
\qty{0.74}{W/m^2}$; increase $2.44 - 0.74 \approx \qty{1.70}{W/m^2}$.

\textbf{12.} One more \unit{W/m^2} multiplies the concentration by
$2^{1/3.93} \approx 1.193$: $427.35 \times 1.193 \approx \qty{510}{ppm}$.

\textbf{13.} $(509.8 - 427.35)/2.634 \approx 31$ years.

\textbf{14.} $2 \times 278 = \qty{556}{ppm}$; \omterm{def:g12:climate-change:forcing}{forcing} $F_{2\times} =
\qty{3.93}{W/m^2}$.

\textbf{15.} By definition of the sensitivity, \qty{3}{\degreeCelsius}
(likely 2.5 to \qty{4}{\degreeCelsius}).

\textbf{16.} $\lambda = 3.93/3 \approx \qty{1.31}{W.m^{-2}.\degreeCelsius^{-1}}$.
The sum of the \omterm{def:g12:climate-change:feedback}{feedbacks}, \qty{1.16}{W.m^{-2}.\degreeCelsius^{-1}}, gives
$3.93/1.16 \approx \qty{3.4}{\degreeCelsius}$, a little more; the assessed
best estimate combines this physics with the observed warming and past
climates, and both values lie in the likely range.

\textbf{17.} The ocean stores most of the extra energy and warms slowly: the
surface reaches its equilibrium warming only after decades to centuries.

\textbf{18.} It assumes one story of future emissions (the growth of the
last ten years kept forever): that is a \omterm{def:g12:climate-change:scenario}{scenario}. A \omterm{def:g12:climate-change:scenario}{projection} computes the
response of the climate and of the carbon cycle to a \omterm{def:g12:climate-change:scenario}{scenario} of emissions
with physical models, and gives ranges and likelihoods.

\textbf{19.} $(556 - 427.35)/2.634 \approx 48.8$ years after 2025: about
\textbf{2074}.
\end{solution}
